\title{
  \iffullversion
    The Quantum Random Oracle Toolbox
    % \thanks{
    %   An extended abstract of this paper will appear at xxx;
    %   this is the full version.
    % }
  \else
    The Quantum Random Oracle Toolbox
    \ifauthversion
    \thanks{
      \copyright IACR 2024. This article is the final version submitted by the authors
      to the IACR and to Springer-Verlag on <date>. The version published
      by Springer-Verlag is available at <DOI>.
    }
    \fi
  \fi
}

\titlerunning{
    The QROT
}

\ifsubmission
  \author{
    \vspace*{-5mm}
  }
  \institute{
    \vspace*{-5mm}
  }
\else
   \author{
    Hans Heum \orcidlink{0000-0003-0527-2999} %\and
  }
  \authorrunning{H.~Heum} 

  \institute{ 
    NTNU - Norwegian University of Science and Technology, Trondheim, Norway.\\
    \email{hans.heum@ntnu.no}
  }
\fi

\maketitle

\begin{abstract}
  Proofs of security in the Quantum Random Oracle Model (QROM) face three main
    challenges, relating to the non-applicability of several conceptually simple yet
    powerful proof techniques from its classical counterpart (ROM):

  \begin{enumerate}
    \item Recording: Or, how to extract QRO queries.
    \item Reprogramming: Or, how to alter specific outputs of the QRO.
    \item Rewinding: Or, how to derandomize quantum oracle algorithms.
  \end{enumerate}

  \noindent
  In this gentle introduction to QROM proof techniques, we will show how to
    deal with the first two. We will
    do this by taking a simple encryption scheme due to Bellare and
    Rogaway~\cite{CCS:BelRog93} as our motivating example, and prove increasingly
    strong statements for it as we move through the literature on quantum
    random oracles and gather stronger and stronger tools.
    %
    We will only assume knowledge of basic linear algebra and standard
    game-hopping techniques (though prior exposure to
    classical random oracles doesn't hurt).

  \begin{keywords}
    %Concrete Security \textperiodcentered~
    Quantum Random Oracles \textperiodcentered~
    One-Way to Hiding \textperiodcentered~
    Compressed Oracle Technique
  \end{keywords}
\end{abstract}

\ifeprintfull
  \newpage
  \tableofcontents
  \newpage
\fi
